\section{Threats to Validity}
\label{sec:threats}

We discuss threats to validity following established guidelines~\cite{wohlin2012experimentation}. Table~\ref{tab:threats-summary} summarizes key threats and mitigations.

\begin{table}[ht]
\centering
\caption{Summary of validity threats and mitigations}
\label{tab:threats-summary}
\begin{tabular}{p{2.2cm}p{4.5cm}}
\toprule
\textbf{Threat} & \textbf{Mitigation} \\
\midrule
LLM non-determinism & 5 runs, statistical analysis, low temperature \\
Prompt confounding & Minimal prompts, P4 control \\
Subject selection & Diverse categories across 11 classes \\
Small sample size & Future work on additional subjects \\
LLM familiarity & P4 (source) as control \\
Manual validation & Two evaluators, inter-rater reliability \\
Mutation limitations & Standard operators, equivalent mutant filtering \\
\bottomrule
\end{tabular}
\end{table}

\paragraph{Internal validity.} LLM non-determinism was mitigated through 5 independent runs with statistical analysis. Prompt confounding was addressed by keeping prompts minimal (differing only in specification inclusion) and adding P4 as a control. Manual validation subjectivity was reduced through clear criteria, two independent evaluators, and high inter-rater reliability (Cohen's $\kappa = 0.91$ for preconditions, $0.87$ for postconditions).

\paragraph{External validity.} We evaluated on Apache Commons Lang (312 methods, 11 classes), which is mature and well-structured, potentially favoring our approach. Results may not generalize to enterprise applications, domain-specific libraries, or proprietary codebases. The LLM may be familiar with Apache Commons Lang from training data, though consistent P1-to-P3 improvement suggests specifications provide value beyond prior knowledge. Our tool and evaluation are Java-specific; extension to other languages requires further research. We used a single LLM (Gemini 2.0); absolute scores may vary with other models, though qualitative findings likely generalize~\cite{schafer2023adaptive, yuan2024chatunitest}.

\paragraph{Construct validity.} Mutation testing has known limitations~\cite{papadakis2019mutation} (equivalent mutants, operator representativeness). We used validated PiTest default operators~\cite{coles2016pitest} and excluded trivially equivalent mutants. Test count is supplemented by mutation score as the primary quality metric.

\paragraph{Conclusion validity.} With 312 methods and 5 runs, power analysis confirms $>$95\% power for observed effect sizes. Bonferroni correction was applied for multiple comparisons. All reported significant differences remain significant after correction.
